\documentclass[11pt]{article}
\usepackage[margin=25mm]{geometry}
\usepackage{fontspec}
\setmainfont{Latin Modern Roman}
\setsansfont{Latin Modern Sans}
\usepackage{amsmath,amssymb,mathtools}
\usepackage{unicode-math}
\setmathfont{Latin Modern Math}
\usepackage{xurl,hyperref}
\setcounter{secnumdepth}{0}
\setlength{\parindent}{0pt}
\setlength{\parskip}{5pt}
\setlength{\emergencystretch}{3em}
\newcommand{\sref}[2]{\hyperlink{source:#1:#2}{[S#2]}}
\newcommand{\cataloguescope}[1]{\paragraph{Recorded target.}#1}
\begin{document}
\section{Kirby Problem 1.64}\label{statement}

\subsection{Definitions and mathematical statement}

\paragraph{Shared notation.}$\mathbb N=\{1,2,\ldots\}$, and $\mathbb Z,\mathbb Q,\mathbb R,\mathbb C$ denote the integers, rationals, reals and complex numbers. Any other conventions are specified locally.


A positively transverse oriented link $L$ in standard contact $(S^3,\xi_{\mathrm{std}})$ is everywhere transverse to the contact planes. The cyclic $n$-fold cover $\Sigma_n(L)$ branched along $L$ has the induced contact structure $\xi_L$: away from the branch set it lifts $\xi_{\mathrm{std}}$, with the usual contact extension near that set. Stein fillability means that this contact manifold is the strictly pseudoconvex boundary of a compact Stein surface. A quasipositive braid is a product of conjugates of positive Artin generators $\sigma_i$ in a braid group; its closure has a natural transverse realization. The self-linking number $\operatorname{sl}(L)$ is the contact push-off linking number computed using a nonzero section over a spanning surface. Slice-Bennequin sharpness means that some smooth compact oriented properly embedded surface $F\subset B^4$, with $\partial F=L$ and no closed components, satisfies $\operatorname{sl}(L)=-\chi(F)$; $\chi$ denotes Euler characteristic.
\paragraph{Statement recorded in the supplied source.}
If $(\Sigma_2(L),\xi_L)$ is Stein fillable, must $L$ be transversely isotopic to the closure of a quasipositive braid? Does the same implication hold if $(\Sigma_n(L),\xi_L)$ is Stein fillable for some integer $n>2$? For a Stein-fillable cyclic branched cover, must the slice-Bennequin inequality for this transverse link be sharp? \sref{statement}{2}

\subsection{Short English statement}
Does Stein fillability of a double or higher cyclic contact branched cover force quasipositivity of the transverse link and sharpness of its slice-Bennequin bound?
\subsection{Sources}
\begin{description}
\item[\hypertarget{source:statement:1}{S1}] ULAM.AI and the UnsolvedMath Contributors, UnsolvedMath: A Curated Collection of Open Mathematics Problems, catalogue entry. CC BY 4.0. \url{https://huggingface.co/datasets/ulamai/UnsolvedMath}.
\item[\hypertarget{source:statement:2}{S2}] R. İnanç Baykur, Robion C. Kirby and Daniel Ruberman, eds., \emph{K3: A New Problem List in Low-Dimensional Topology}, Mathematical Surveys and Monographs 295, American Mathematical Society (2026), Problem 1.64 and its remarks; Chapters 1 and 2 for the stated conventions. \url{https://math.berkeley.edu/sites/default/files/surv-295-ruberman-watermarked-author-pdf.pdf}.
\end{description}
\label{end:statement}
\end{document}
